\documentclass{rrxiv}
\rrxivid{rrxiv:2605.00006}
\rrxivversion{v1}
\rrxivprotocolversion{0.1.0}
\rrxivlicense{CC-BY-4.0}
\rrxivtopics{cs.DL}

\title{Citation graphs are not knowledge graphs}
\author{Example Author A}
\date{2026-05-17}

\begin{document}
\maketitle

\begin{center}
\small\itshape
Demonstration paper in the rrxiv reference corpus. The canonical machine-readable version lives at \href{https://rrxiv.com/papers/rrxiv:2605.00006}{rrxiv.com/papers/rrxiv:2605.00006}.
\end{center}

\begin{abstract}
An opinion piece distinguishing the structural commitments of citation networks from those of knowledge graphs proper, with implications for preprint infrastructure. We argue that conflating the two has produced a decade of misnamed retrieval systems.
\end{abstract}

\section{Introduction}
An opinion piece distinguishing the structural commitments of citation networks from those of knowledge graphs proper, with implications for preprint infrastructure. We argue that conflating the two has produced a decade of misnamed retrieval systems.

This document is a structured encoding of the paper in the \texttt{rrxiv} protocol's Canonical Intermediate Representation (CIR). It engages with the topic \texttt{cs.DL}. The encoding registers 1 formal claim (1 untested). Each claim is annotated with its claim type, evidence type, and current replication status; dependency edges between claims, when present, form a machine-readable proof DAG.

\section{Methodology}
We follow the \texttt{rrxiv} convention of separating \emph{claims} (the proposition under consideration) from \emph{evidence} (the argument or data supporting it). Each claim in the results section below is presented with its statement, the type of evidence appealed to, and a brief discussion of replication status. Where claims depend on prior results --- internal or external --- the dependency is recorded in the CIR as a \texttt{\textbackslash dependson} edge, so the full inferential structure is machine-traversable. Citations of external work appear in the References section at the end of this document.

\section{Results: registered claims}
\subsection*{Claim 1}
\begin{claim}[Claim 1]
\label{claim:c1}
Citation networks encode 'this work cited that work' edges; knowledge graphs encode typed semantic relations — the two are not interchangeable.

\emph{Replication status: untested.}
\end{claim}
As of the encoding date, it has not yet been independently tested. It depends on 1 prior claim in the same paper.

\section{Discussion}
The claim graph above is the primary product of this paper. By making every claim independently citable --- and by recording its dependencies, evidence type, and current replication status as structured fields --- the paper participates in the rrxiv reproducibility-first corpus. Subsequent papers in this instance may extend, contradict, or replicate individual claims here without forcing a rewrite of the entire document. See the canonical version online for the live discourse layer.

\section{References}
\begin{itemize}[leftmargin=*]
\item Knowledge graphs from scientific abstracts
\item Citation networks vs knowledge graphs
\end{itemize}
\end{document}
